\documentclass[12pt]{article}
% Préambule commun à tous les documents extraits de « La Revue CAPES »
\usepackage[T1]{fontenc}
\usepackage[utf8]{inputenc}
\usepackage{lmodern}
\usepackage{amsmath,amssymb,amsthm,amsfonts,mathrsfs}
\usepackage{setspace}
\usepackage{listings}
\usepackage{pgfplots}
\pgfplotsset{compat=1.18}
\usepackage{longtable}
\usepackage{adjustbox}
\usepackage{lettrine}
\usepackage{tkz-tab}
\usepackage{changepage}
\usepackage{pdflscape}
\usepackage{graphicx}
\usepackage{tikz}
\usepackage{xcolor}
\usepackage[a4paper,top=2cm,bottom=2.5cm,left=2.5cm,right=2cm]{geometry}
\usepackage{fancyhdr}
\usepackage{draftwatermark}
\SetWatermarkText{La RevueCapes}
\SetWatermarkLightness{0.9}
\SetWatermarkScale{2}
\usepackage{hyperref}
\hypersetup{colorlinks=true,linkcolor=blue!50!black,urlcolor=blue!50!black}

\definecolor{revue}{RGB}{20,60,120}

\pagestyle{fancy}
\fancyhf{}
\renewcommand{\headrulewidth}{0.4pt}
\setlength{\headheight}{15pt}

% \entete{Matière}{Titre}{Type de document}
\newcommand{\entete}[3]{%
  \fancyhead[L]{\small\textsc{La Revue CAPES} -- #1}%
  \fancyhead[R]{\small #2}%
  \fancyfoot[C]{\thepage}%
  \thispagestyle{plain}%
  \begin{center}
    {\color{revue}\rule{\textwidth}{1.2pt}}\\[4pt]
    {\small\textsc{Concours directs CAP-CEG / CAPES -- Burkina Faso}}\\[6pt]
    {\LARGE\bfseries\color{revue} #2}\\[6pt]
    {\large\itshape #1 -- #3}\\[2pt]
    {\color{revue}\rule{\textwidth}{1.2pt}}
  \end{center}
  \vspace{0.5cm}%
}

% Encadré pour les remarques ajoutées dans les corrigés
\newcommand{\remarque}[1]{\par\medskip\noindent\fcolorbox{revue}{revue!6}{\parbox{\dimexpr\linewidth-2\fboxsep-2\fboxrule}{\textbf{\color{revue}Remarque.} #1}}\par\medskip}


\begin{document}
\entete{Mathématiques}{Corrigé détaillé du sujet type n°10}{Correction}
\begin{spacing}{1.5}

\textbf{Exercice 1}

$\begin{cases}
ax-by=c\\
x-2y=3
\end{cases}$

1) Il y a $6^3=216$ possibilités pour le triplets $(a,b,c)$.

2) Calculons la probabilité $P_1$ que $(S)$ ait une infinité de solutions.

Pour que $(S)$ ait une infinité de solution, les deux équations doivent être proportionnelles, soit : 

$a = 1, b = 2$ et $c = 3$ ou 

$a = 2, b = 4$ et $c = 6$  

Aucune autre configuration de tirages ne conduit à la même équation $x-2y = 3$. 

La probabilité $P_1$ pour que $(S)$ ait une infinité de solutions est donc $2/216 = 1/108$. 

3) Calculons la probabilité $P_2$ que $(S)$ n'admette aucune solution. 

Pour qu'il n'y ait pas de solution, il faut que $b-2a = 0$. 

Cela conduit à : 

$a = 1, b = 2$, et $c\neq 3$ (car si $c = 3$, on est dans le cas précédent) 

$a = 2, b = 4$, et $c\neq 6$ 

$a = 3, b = 6$, et $c$ prend n'importe quelle valeur. 

Soit au total 5 + 5 + 6 = 16 cas possibles. 

La probabilité $P_2$ pour que $(S)$ n'admette aucune solution est 16/216 = 8/108. 

4) Calculons la probabilité $P_3$ que $(S)$ admette une solution unique. 

Il existe une et une seule solution si $b-2a\neq 0$. 

Soit : 
$a = 1, b = 1, 3, 4, 5, 6 $, $\forall c$
 
a$ = 2, b = 1, 2, 3, 5, 6 $, $\forall c $

$a = 2, b = 1, 2, 3, 4, 5 $, $\forall c $

$a = 4$, $\forall b$, $\forall c$ 

$a = 5$, $\forall b$, $\forall c$
 
$a = 6$, $\forall b$, $\forall c$
 
Cela représente 3 x 6 x 5 + 3 x 36 = 198 cas possibles. 

La probabilité $P_3$ pour que $(S)$ admette une solution unique est 198/218 = 99/108. 

On vérifie $P_1 + P_2 + P_3 = 1$.

5) Calculons la probabilité $P_4$ que $(S)$ admette le couple $(x = 3, y = 0)$ comme unique 
solution.

Si (3,0) est solution : $3a-0 = c 
\Rightarrow a = 1$ et $c = 3$, ou $a = 2$ et $c = 6$, et $b$ est quelconque. 

$P_4 = 12/216 = 6/108$

\medskip\textbf{Exercice 2}


On considère la série entière:

\[\sum_{n=1}^{+\infty}\frac{x^n}{n(n+2)}\]

1) Déterminons son rayon de convergence $R$.

Soit $u_n=\frac{1}{n(n+2)}$ donc $u_{n+1}=\frac{1}{(n+1)(n+3)}$

\[R=\frac{1}{l} \text{ avec } l=\lim_{n\to +\infty}\frac{u_{n+1}}{u_n}=\frac{n(n+2)}{(n+1)(n+3)}=\lim_{n\to +\infty}\frac{n^2}{n^2}=1\\
\text{ ainsi } R=1\]

2) Étudions la nature de la série pour $x=1$ et $x=-1$.

$-$ Si $x=1$, on a :

\[\sum_{n=1}^{+\infty}\frac{1}{n(n+2)}\sim\sum_{n=1}^{+\infty}\frac{1}{n^2}\]

$\sum_{n=1}^{+\infty}\frac{1}{n^2}$ est Riemann convergente, alors $\sum_{n=1}^{+\infty}\frac{1}{n(n+2)}$ converge.

\quad
$-$ Si $x=-1$, on a $\sum_{n=1}^{+\infty}\frac{(-1)^n}{n(n+2)}$

Posons $u_n=\frac{1}{n(n+2)}$

$\lim_{n\to +\infty}u_n=\lim_{n\to +\infty}\frac{1}{n(n+2)}=0$

$\frac{u_{n+1}}{u_n}=\frac{n(n+2)}{(n+1)(n+3)}< 1$ (puisque $n<n+1$ et $n+2<n+3$). Alors $u_{n+1}<u_n$, d'où $u_n$ est décroissante.

Comme $\lim_{n\to +\infty}u_n=0$ et $u_n$ est décroissante, alors on conclut que $u_n$ est convergente d'après le critère des série alternées



3) Décomposons en éléments simples la fraction rationnelle $\frac{1}{X(X+2)}$.

$\frac{1}{X(X+2)}=\frac{a}{X}+\frac{b}{X+2}$

En multipliant chaque terme par $x$, et en prenant $X=0$, on obtient $a=\frac{1}{2}$

En multipliant chaque terme par $x+2$, et en prenant $X=-2$, on obtient $b=-\frac{1}{2}$

Alors $\frac{1}{X(X+2)}=\frac{1}{2}(\frac{1}{X}-\frac{1}{X+2})$

4) Rappelons le développement en série entière de la fonction x$\longmapsto ln(1-x)$ en 0.

On a :$\sum_{n=0}^{+\infty}x^n=-\frac{1}{1-x}$

En appliquant l'intégrale , on obtient :$\sum_{n=0}^{+\infty}\frac{x^{n+1}}{n+1}=-\ln(1-x)$

Posons : \[p=n+1\Rightarrow \sum_{p=1}^{+\infty}\frac{x^p}{p}=-ln(1-x)\Rightarrow \ln(1-x)= -\sum_{p=1}^{+\infty}\frac{x^p}{p}=-\sum_{n=1}^{+\infty}\frac{x^n}{n}\]

5) $\forall x\in ]-R,R[$, Déduisons la somme:

$S=\sum_{n=1}^{+\infty}\frac{x^n}{n+2}$.

\[S=\sum_{n=1}^{+\infty}\frac{x^n}{n+2}=\sum_{n=1}^{+\infty}\frac{x^{n+2-2}}{n+2}=\sum_{n=1}^{+\infty}\frac{x^{n+2}}{n+2}\times\frac{1}{x^2}=\frac{1}{x^2}\sum_{n=1}^{+\infty}\frac{x^{n+2}}{n+2}\]

Posons : $p=n+2$, donc $S=\frac{1}{x^2}\sum_{p=3}^{+\infty}\frac{x^{p}}{p}$

$p$ est une variable muette, donc $S=\frac{1}{x^2}\sum_{p=3}^{+\infty}\frac{x^{p}}{p}=\frac{1}{x^2}\sum_{n=3}^{+\infty}\frac{x^{n}}{n}$


On a : $\ln(1-x)=-\sum_{n=1}^{+\infty}\frac{x^n}{n}=-\left( x+\frac{x^2}{2}+\sum_{n=3}^{+\infty}\frac{x^n}{n}\right)$ d'où  $\sum_{n=3}^{+\infty}\frac{x^n}{n}=-\ln(1-x)-x-\frac{x^2}{2}$

Donc $S=\frac{1}{x^2}\left(\ln(1-x)-x-\frac{x^2}{2}\right)=\frac{-2\ln(1-x)-2x-x^2}{2x^2}$


6) Montrons que $\forall x\in ]-R,0[\cup ]0,R[$, $\sum_{n=1}^{+\infty}\frac{x^n}{n(n+2)}=\frac{2(1-x^2)ln(1-x)+2x+x^2}{4x^2}$.


\begin{align*}
\sum_{n=1}^{+\infty}\frac{x^n}{n(n+2)} &=\sum_{n=1}^{+\infty}\frac{1}{2}\left(\frac{1}{n}-\frac{1}{n+2}\right)x^n =\frac{1}{2}\sum_{n=1}^{+\infty}\frac{x^n}{n} - \frac{1}{2}\sum_{n=1}^{+\infty}\frac{x^n}{n+2}\\
&=\frac{1}{2}(-\ln(1-x))-\frac{1}{2}\frac{-2\ln(1-x)-2x-x^2}{2x^2}\\
&=\frac{1}{2}\frac{-2x^2\ln(1-x)+2\ln(1-x)+2x+x^2}{2x^2}\\
 &=\frac{2(1-x^2)\ln(1-x)+2x+x^2}{4x^2}\text{  } \forall x\in ]-R,0[\cup ]0,R[
\end{align*}


7) Via un équivalent de $\ln(1-x)$, vérifions que l'on peut prolonger cette dernière expression en 0.


En $x=0$, $\ln(1-x)\sim x$, alors :

\[\lim_{x\to 0}\frac{2(1-x^2)\ln(1-x)+2x+x^2}{4x^2}=\lim_{x\to 0}\frac{2(1-x^2)x+2x+x^2}{4x^2}=\lim_{x\to 0}\frac{2x+1}{4}=\frac{1}{4}\]

8) En utilisant la question précédentes calculons les sommes suivantes :

\[S_1=\sum_{n=2}^{+\infty}\frac{(-1)^n}{n(n+2)} \text{ et } S_2=\sum_{n=2}^{+\infty}\frac{1}{n(n+2)}\]

On a :

\[* \sum_{n=1}^{+\infty}\frac{(-1)^n}{n(n+2)}=\frac{2(1-(-1)^2)\ln(1-(-1))+2\times(-1)+(-1)^2}{4(-1)^2}=-\frac{1}{4}\]
\[\sum_{n=1}^{+\infty}\frac{(-1)^n}{n(n+2)}=-\frac{1}{3}+\sum_{n=2}^{+\infty}\frac{(-1)^n}{n(n+2)}=-\frac{1}{3}+S_1\Rightarrow S_1=\frac{1}{3}-\frac{1}{4}=\frac{1}{12}\]


\[* \sum_{n=1}^{+\infty}\frac{1}{n(n+2)}=\frac{3}{4}=\frac{1}{3}+\sum_{n=2}^{+\infty}\frac{1}{n(n+2)}=\frac{1}{3}+S_2\Rightarrow S_2 =\frac{3}{4}-\frac{1}{3}=\frac{5}{12}\]







\medskip\textbf{Exercice 3}

1. Démonstration des valeurs propres de $A$

Les valeurs propres de $A$ sont les solutions de l'équation caractéristique $\det(A - \lambda I) = 0$, où $I$ est la matrice identité. $
A - \lambda I = 
\begin{pmatrix}
1-\lambda & 0 & 1 \\ 
-1 & 2-\lambda & 1 \\ 
1 & -1 & 1-\lambda
\end{pmatrix}.
$

Le déterminant est donné par :
$
\det(A - \lambda I) = \begin{vmatrix}
1-\lambda & 0 & 1 \\ 
-1 & 2-\lambda & 1 \\ 
1 & -1 & 1-\lambda
\end{vmatrix}.
$

En développant suivant la première ligne :
\[
\det(A - \lambda I) = (1-\lambda) \begin{vmatrix}
2-\lambda & 1 \\ 
-1 & 1-\lambda
\end{vmatrix} 
+ \begin{vmatrix}
-1 & 2-\lambda \\ 
1 & -1
\end{vmatrix}.
\]

Les mineurs sont :
$
\begin{vmatrix}
2-\lambda & 1 \\ 
-1 & 1-\lambda
\end{vmatrix} = \lambda^2 - 3\lambda + 3,$ et 
$
\begin{vmatrix}
-1 & 2-\lambda \\ 
1 & -1
\end{vmatrix} = -1+\lambda.
$

Ainsi :
\[
\det(A - \lambda I) = (1-\lambda)(\lambda^2 - 3\lambda + 3) +\lambda -1=
\]

En factorisant, on obtient :
\[
\det(A - \lambda I) =(\lambda-1)(-\lambda^2 + 3\lambda - 3+1)=(\lambda-1)(-\lambda^2 + 3\lambda -2)= (\lambda - 1)^2 (\lambda - 2).
\]

Les valeurs propres sont donc :
$
\lambda = 1 \quad (\text{multiplicité 2}) \quad \text{et} \quad \lambda = 2 \quad (\text{multiplicité 1}).
$

2. Sous-espaces propres et diagonalisabilité de $A$

Valeur propre $\lambda = 1$

La matrice $A - I$ est :
$
A - I = 
\begin{pmatrix}
0 & 0 & 1 \\ 
-1 & 1 & 1 \\ 
1 & -1 & 0
\end{pmatrix}.
$

Résolvons $(A - I)X = 0$. Le système est :
$
\begin{cases}
x + z = 0, \\
-y + z = 0, \\
x - y = 0.
\end{cases}
$

On trouve une base de $E_1$ :
$
E_1 = \mathrm{Vect}\left(\begin{pmatrix} 1 \\ 1 \\ -1 \end{pmatrix}\right).
$

Valeur propre $\lambda = 2$

La matrice $A - 2I$ est :
$
A - 2I = 
\begin{pmatrix}
-1 & 0 & 1 \\ 
-1 & 0 & 1 \\ 
1 & -1 & -1
\end{pmatrix}.
$

Résolvons $(A - 2I)X = 0$. Le système est :
$
\begin{cases}
-x + z = 0, \\
-y + z = 0, \\
x - y = 0.
\end{cases}
$

On trouve une base de $E_2$ :
$
E_2 = \mathrm{Vect}\left(\begin{pmatrix} 1 \\ 1 \\ 1 \end{pmatrix}\right).
$

Diagonalisabilité

La matrice $A$ n'est \textbf{pas diagonalisable}, car la dimension de $E_1$ est strictement inférieure à la multiplicité algébrique de $\lambda = 1$.

3. Sous-espaces caractéristiques

Le sous-espace caractéristique de $\lambda = 1$ est donné par $F_1 = \ker((A - I)^2)$.

Calculons :
\[
(A - I)^2 = 
\begin{pmatrix}
0 & 0 & 1 \\ 
-1 & 1 & 1 \\ 
1 & -1 & 0
\end{pmatrix}^2 = 
\begin{pmatrix}
1 & -1 & 0 \\ 
-1 & 1 & 0 \\ 
0 & 0 & 0
\end{pmatrix}.
\]

Résolvons $(A - I)^2X = 0$. On trouve :
$
F_1 = \mathrm{Vect}\left(\begin{pmatrix} 1 \\ 1 \\ -1 \end{pmatrix}, \begin{pmatrix} 1 \\ 0 \\ 0 \end{pmatrix}\right).
$

Le sous-espace caractéristique pour $\lambda = 2$ est $F_2 = E_2$.

4. Base pour obtenir $B$ et décomposition de Dunford

Une base permettant de transformer $A$ en une matrice $B$ est :
\[
\mathcal{B} = \left\{\begin{pmatrix} 1 \\ 1 \\ 1 \end{pmatrix}, \begin{pmatrix} 1 \\ 1 \\ -1 \end{pmatrix}, \begin{pmatrix} 1 \\ 0 \\ 0 \end{pmatrix}\right\}.
\]

Dans cette base, $A$ devient :
$
B = \begin{pmatrix}
2 & 0 & 0 \\ 
0 & 1 & 1 \\ 
0 & 0 & 1
\end{pmatrix}.
$

La décomposition de Dunford est donnée par :$
B = D + N,
$
avec :
$
D = \begin{pmatrix}
2 & 0 & 0 \\ 
0 & 1 & 0 \\ 
0 & 0 & 1
\end{pmatrix}, \quad
N = \begin{pmatrix}
0 & 0 & 0 \\ 
0 & 0 & 1 \\ 
0 & 0 & 0
\end{pmatrix}.
$

5. Résolution du système différentiel

Le système est :
$
\begin{cases}
x' = x + z, \\
y' = -x + 2y + z, \\
z' = x - y + z.
\end{cases}
$

En notation matricielle :
\[
X'(t) = AX(t),
\]
où $A = \begin{pmatrix}
1 & 0 & 1 \\ 
-1 & 2 & 1 \\ 
1 & -1 & 1
\end{pmatrix}.
$

La solution est donnée par :$
X(t) = e^{tA}X(0).
$

En utilisant $A = D + N$, on trouve :$
e^{tA} = e^{tD}e^{tN}.
$

Avec :
$
e^{tD} = \begin{pmatrix}
e^{2t} & 0 & 0 \\ 
0 & e^t & 0 \\ 
0 & 0 & e^t
\end{pmatrix}, \quad
e^{tN} = \begin{pmatrix}
1 & 0 & 0 \\ 
0 & 1 & t \\ 
0 & 0 & 1
\end{pmatrix}.
$ 
Ainsi :
$
e^{tA} = \begin{pmatrix}
e^{2t} & 0 & 0 \\ 
0 & e^t & te^t \\ 
0 & 0 & e^t
\end{pmatrix}.
$

La solution est :
$
X(t) = \begin{pmatrix}
e^{2t}x_0 \\ 
e^ty_0 + te^tz_0 \\ 
e^tz_0
\end{pmatrix}.
$


\medskip\textbf{Exercice 4}

Trouvons toutes les applications $f :\mathbb{R} \longrightarrow \mathbb{R}$ deux fois dérivables telles que: 
$\forall x \in \mathbb{R}$ , $f"(x) + f(-x) = x$ $(E)$.

L'équation proposée $(E)$ n'est pas une équation différentielle.
En remplaçant $x$ par $-x$ dans $(E)$, on obtient : $f"(-x) + f(x) = -x$

D'où le système :

\[
\begin{cases}
f"(x) + f(-x) = x\\
f"(-x) + f(x) = -x
\end{cases}
\]

En additionnant membre par membre ces deux équations, puis en les soustrayant, on obtient en tenant compte des notations des fonctions $f$ et $g$ : 

\[
\begin{cases}
[f"(x)+f"(-x)]+[f(x)+f(-x)] =g"(x) + g(x) =0\\
[f"(x)-f"(-x)]+[f(x)-f(-x)] =h"(x)-h(x)=2x
\end{cases}
\]

La première équation admet pour solution

 $g(x) = C_1\cos x + C_2\sin x$ avec $(C_1,C_2 )\in \mathbb{R}^2$

La seconde équation admet pour solution $h(x) =-2x + D_1chx + D_2shx$ avec $(D_1,D_2 )\in\mathbb{R}^2$ et donc puisque

$f(x) =\frac{1}{2}[g(x) +h(x)]=-x +\frac{1}{2}[C_1 \cos x + C_2 \sin x + D_1chx + D_2shx]$

Étudions la réciproque :

Les solutions $f(x)$ vérifiant l'équation $(E)$ 
$f"(x) + f(-x) = 0$ $\forall x\in \mathbb{R}$ doivent satisfaire
 
$C_2\sin x + D_1chx = 0$  $\forall x\in \mathbb{R}$ et donc $C_2 = D_1 = 0$

Les solutions de $(E)$ sont donc : $f(x) =-x+\frac{1}{2}[C_1\cos x + D_2shx]$ avec $(C_1,D_2)\in\mathbb{R}^2$

\medskip\underline{Remarque} :

Par construction, $g$ est une fonction paire ce qui implique $C_2 = 0$ et $h$ une fonction impaire ce
qui implique $D_1 =0$ 

\end{spacing}
\end{document}
